\question[5]
Gib für jede Anweisung an, ob eine implizite oder eine explizite Typumwandlung stattfindet, welchen Wert die
Variable danach hat und welches Problem dabei auftritt (falls eines auftritt). Bei \texttt{e} genügt es, die
Größenordnung des Wertes abzuschätzen.

Zur Erinnerung: Ein \texttt{float} speichert höchstens etwa 16 Dezimalstellen.

\begin{lstlisting}
a = 3 + 0.5
b = int(9.99)
c = int('3.5')
d = bool('False')
e = float(2 ** 60 + 1)
\end{lstlisting}

\begin{solutionbox}{6cm}
\mbox{}\\
\begin{tabular}[t]{lllp{7.5cm}}
 & Umwandlung & Wert & Problem \\ \hline
\texttt{a} & implizit & \texttt{3.5} & keines, \texttt{int} wird zu \texttt{float} \\
\texttt{b} & explizit & \texttt{9} & Informationsverlust: Nachkommastellen werden abgeschnitten, nicht gerundet \\
\texttt{c} & explizit & -- & \texttt{ValueError}: \texttt{'3.5'} ist keine ganze Zahl \\
\texttt{d} & explizit & \texttt{True} & überraschend: jeder nicht leere String ergibt \texttt{True} \\
\texttt{e} & explizit & ca. $10^{18}$ & Genauigkeitsverlust: $2^{60} = (2^{10})^6 \approx (10^3)^6 = 10^{18}$
hat 19 Stellen, ein \texttt{float} speichert aber höchstens etwa 16. Die letzten Stellen gehen verloren,
die $+1$ verschwindet. (Genauer Wert: \texttt{1.152921504606847e+18}) \\
\end{tabular}
\end{solutionbox}
